% ============================================
% ===== CAPÍTULO 1: PLANTEAMIENTO DEL PROBLEMA =====
% ============================================

\chapter{Planteamiento del problema}
\label{chap:planteamiento}

\section{Contexto}
\label{sec:contexto}

Spotify es una de las plataformas de streaming musical con mayor número de usuarios a nivel mundial. A través de su API pública (\emph{Spotify Web API}) permite, previa autenticación del propio usuario, consultar información agregada sobre sus hábitos de escucha: artistas más reproducidos, canciones favoritas, listas de reproducción, etc.

Este proyecto nace con el objetivo de aplicar el ciclo de vida de un proyecto de ciencia de datos ---presentado en la primera sesión de la asignatura--- sobre un conjunto de datos propio y personal: el historial de artistas más escuchados de los autores de este trabajo, enriquecido con información de género y de similitud obtenida de Last.fm y con una segunda medida de popularidad procedente de Deezer.

\section{Análisis de negocio}
\label{sec:analisis-negocio}

Siguiendo el esquema de preguntas planteado en la asignatura para delimitar el objetivo de un proyecto de ciencia de datos (¿cuánto?, ¿qué categoría?, ¿qué grupo?, ¿es esto raro?, ¿qué opción se debe tomar?), este trabajo se sitúa principalmente en torno a las preguntas de \textbf{categorización}, \textbf{agrupamiento} y \textbf{comparación}:

\begin{itemize}
    \item ¿Qué géneros musicales predominan en los gustos de cada usuario?
    \item ¿Es estable el gusto musical en el tiempo, o varía al comparar el corto plazo (\emph{short\_term}, últimas 4 semanas), el medio plazo (\emph{medium\_term}, últimos 6 meses) y el largo plazo (\emph{long\_term}, varios años)?
    \item ¿Existen agrupaciones naturales de artistas que permitan describir el perfil musical de forma más compacta que una simple lista?
    \item ¿Qué diferencias y similitudes existen entre los perfiles musicales de los dos autores del trabajo?
    \item ¿Se comportan los autores de forma parecida respecto a la popularidad de los artistas que escuchan?
    \item ¿Se puede estimar, para un artista ajeno al historial, si encaja con el gusto de un usuario a partir de la similitud con sus artistas favoritos?
\end{itemize}

\section{Objetivo del proyecto}
\label{sec:objetivo}

El objetivo principal es construir, a partir de datos personales obtenidos mediante la API de Spotify, un conjunto de datos estructurado que recoja los artistas más escuchados por cada autor en los tres rangos temporales disponibles, junto con los géneros musicales asociados a cada artista y otros atributos de popularidad y similitud. Sobre ese conjunto de datos se realiza:

\begin{itemize}
    \item Un análisis descriptivo de la diversidad y la evolución del gusto musical (distribución de géneros, comparación entre rangos temporales y entre usuarios).
    \item Un análisis de compatibilidad entre los dos autores mediante métricas de similitud sobre artistas y géneros.
    \item Un análisis de popularidad, relacionando el ranking personal con los oyentes de Last.fm y contrastando estos con los fans de Deezer.
    \item Una estimación de afinidad usuario--artista para artistas nuevos, basada en la vecindad de similitud proporcionada por Last.fm.
    \item Un \emph{clustering} exploratorio de artistas por sus etiquetas, incorporado como ampliación metodológica propia del proyecto.
    \item La visualización de los resultados mediante los tipos de gráficos estudiados en la asignatura, justificando en cada caso el tipo elegido.
\end{itemize}

\section{Marco metodológico y alcance}

La organización sigue las fases de comprensión del problema, comprensión y preparación de los datos, modelado, evaluación y comunicación de CRISP-DM, presentada en la Sesión~1~\cite{sesion1}. La entrega es un estudio exploratorio, no el despliegue de un sistema de recomendación. La selección de fuentes, limpieza y formatos se fundamenta en las Sesiones~1--3, y la comunicación visual en la Sesión~6. Jaccard, coseno, Spearman y K-Means son elecciones metodológicas del proyecto, no requisitos atribuidos al material docente.

Los CSV contienen \textbf{rankings agregados}, no un registro de reproducciones: no se conoce el tiempo de escucha ni cuántas veces se ha reproducido cada artista. Los rangos temporales se solapan y no son tres observaciones longitudinales independientes. Por ello se describe la coincidencia entre ventanas, sin inferir tendencias causales ni generalizar a una población.

El enunciado permite elegir las fuentes y el problema~\cite{enunciado}; no exige usar todas las técnicas de las sesiones. No se añaden análisis geográficos ni de grafos porque las preguntas delimitadas pueden abordarse con las variables disponibles.
